\chapter{The Pro-etale--Classical Adic Dictionary}
\label{app:proetale-classical-dictionary}

The pro-etale category does not define a new Weil cohomology theory.  It gives
a direct model for the classical \(\ell\)-adic category.  This appendix states
the comparison in the precise form used by the notes and records which
expressions correspond only by canonical isomorphism rather than by literal
equality.

\section{The two integral models}

Let \(X\) be topologically noetherian, let
\(\mathcal O\) be a complete discrete valuation ring with uniformizer
\(\varpi\), and put \(\Lambda_n=\mathcal O/\varpi^n\).

The classical model begins with projective systems
\[
  K_\bullet=(K_n)_{n\geq1},
  \qquad
  K_n\in\Dbc(X_{\et},\Lambda_n),
\]
equipped with derived normalization isomorphisms
\[
  K_{n+1}\otimes^L_{\Lambda_{n+1}}\Lambda_n
     \xrightarrow{\sim}K_n.
\tag{B.1}\label{eq:dictionary-normalized-system}
\]
One then localizes an appropriate derived category of systems at AR-null
complexes.  We denote the result schematically by
\(\mathrm D^b_{c,\mathrm{adic}}(X,\mathcal O)\).

The pro-etale model uses the actual sheaf
\[
  \widehat{\mathcal O}_X
     =\varprojlim_n(\Lambda_n)_X
\]
and the full category
\[
  \Dcons(X_{\proet},\widehat{\mathcal O}_X)
  \subset
  \D(X_{\proet},\widehat{\mathcal O}_X)
\]
of derived-complete constructible complexes.

\begin{theorem}[Integral comparison dictionary]
\label{thm:dictionary-integral-comparison}
There is a canonical equivalence
\[
 \Phi_X:
 \mathrm D^b_{c,\mathrm{adic}}(X,\mathcal O)
 \xrightarrow{\sim}
 \Dcons(X_{\proet},\widehat{\mathcal O}_X)
\tag{B.2}\label{eq:dictionary-integral-equivalence}
\]
characterized on a normalized system by
\[
 \Phi_X(K_\bullet)
   \simeq
 \Rlim_n\nu^*K_n,
\tag{B.3}\label{eq:dictionary-realization}
\]
where \(\nu:X_{\proet}\to X_{\et}\).  Its inverse sends \(K\) to the
compatible reductions
\[
 K_n=K\otimes^L_{\widehat{\mathcal O}_X}\Lambda_n.
\tag{B.4}\label{eq:dictionary-reductions}
\]
After inverting \(\varpi\), it gives
\[
 \mathrm D^b_{c,\mathrm{adic}}(X,E)
 \xrightarrow{\sim}
 \Dcons(X_{\proet},E_X).
\tag{B.5}\label{eq:dictionary-rational-equivalence}
\]
\end{theorem}

\begin{proof}
Derived completeness gives
\[
 K\simeq\Rlim_n K_n.
\]
Every \(K_n\) is classical by the torsion comparison
\(\D^+(X_{\et},\Lambda_n)\to\D^+(X_{\proet},\Lambda_n)\).
Conversely a normalized constructible system has a derived inverse limit;
repleteness and the uniform Tor-amplitude bound show that reduction of this
limit is \(K_n\).  A morphism is an isomorphism exactly when all finite
reductions are isomorphisms.  Hence the two constructions are inverse after
the classical localization by pro-zero errors.  Inverting \(\varpi\) gives
\eqref{eq:dictionary-rational-equivalence}.  This is the content of the
comparison developed in Bhatt--Scholze Sections 6.5--6.8.
\end{proof}

\begin{proetalebox}
The equivalence explains the simplification.  AR-null complexes have not
been declared nonexistent; their localization has been absorbed into
derived completion in a replete topos.  The public object is \(K\), and its
coefficient tower is recovered functorially as \((K/\varpi^n)_n\).
\end{proetalebox}

\section{Object-by-object dictionary}

\begin{longtable}{
 >{\raggedright\arraybackslash}p{0.23\textwidth}
 >{\raggedright\arraybackslash}p{0.31\textwidth}
 >{\raggedright\arraybackslash}p{0.36\textwidth}}
\toprule
Concept & Classical adic model & Pro-etale model\\
\midrule
\endhead
Coefficient ring &
The tower \((\Lambda_n)_n\). &
The sheaf
\(\widehat{\mathcal O}_X=\varprojlim_n(\Lambda_n)_X\).\\

Integral complex &
Normalized constructible system, modulo AR-null systems. &
Derived-complete \(K\) with constructible perfect reduction.\\

Reduction modulo \(\varpi^n\) &
The \(n\)-th normalized component \(K_n\). &
\(K\otimes^L_{\widehat{\mathcal O}}\Lambda_n\).\\

Completion &
Renormalization of the coefficient tower. &
The idempotent functor
\[
 K^\wedge=\Rlim_n(K\otimes^L\Lambda_n).
\]\\

Lisse integral object &
A compatible tower of finite locally constant
\(\Lambda_n\)-modules, locally finite free. &
A locally finite-free
\(\widehat{\mathcal O}\)-sheaf; derived lisse objects are locally perfect.\\

Constructible object &
Lisse normalized systems along a finite stratification. &
Locally perfect derived objects along a finite stratification.\\

Rational object &
An integral lattice with morphisms tensored by \(E\), up to commensurability. &
An actual constructible \(E_X\)-complex; equivalently an integral
constructible object after inverting \(\varpi\).\\

\(\Qbarl\)-object &
Extension from a sufficiently large finite coefficient field. &
The same finite-field descent, expressed using the sheaf associated with the
topological algebraic extension.\\

Stalk &
The compatible system \((K_{n,\bar x})_n\), then inverse limit. &
The pullback of the actual complete complex to the geometric point.\\

Cohomology &
\(\Rlim_nR\Gamma(X_{\et},K_n)\). &
\(R\Gamma(X_{\proet},K)\) literally.\\

Continuous representation &
A compatible system of finite continuous representations of the classical
profinite etale fundamental group. &
For connected locally noetherian \(X\) with a geometric base point \(\bar x\),
a locally constant pro-etale sheaf, equivalently a continuous representation
of the Noohi group \(\pi_1^{\proet}(X,\bar x)\).\\
\bottomrule
\end{longtable}

The base-pointed representation statement in the last row is made only under
the displayed connectedness and local-noetherian hypotheses.  Finite
discrete representations, and those \(E\)-representations that stabilize an
\(\mathcal O_E\)-lattice, factor through the classical profinite
\(\pi_1^{\et}(X,\bar x)\).  On geometrically unibranch schemes this captures
the lisse \(E\)-objects relevant to the proof.  For a general nonnormal
scheme, however, the Noohi group \(\pi_1^{\proet}\) can be strictly larger,
and a continuous finite-dimensional \(E\)-representation need not preserve a
lattice.  The Noohi and classical statements must therefore not be identified
globally.

\section{The six-functor dictionary}

Write \(K_n=K\otimes^L\Lambda_n\).  Under
\eqref{eq:dictionary-integral-equivalence}, every operation is characterized
by finite reduction and its adjunction.

\begin{longtable}{
 >{\raggedright\arraybackslash}p{0.19\textwidth}
 >{\raggedright\arraybackslash}p{0.34\textwidth}
 >{\raggedright\arraybackslash}p{0.37\textwidth}}
\toprule
Operation & Pro-etale construction & Comparison formula\\
\midrule
\endhead
Pullback &
Completed derived pullback
\(Lf_{\comp}^*K=(Lf^*K)^\wedge\). &
\[
 (Lf_{\comp}^*K)/\varpi^n
   \simeq Lf^*(K_n).
\]\\

Direct image &
Derived pushforward \(Rf_*\), which preserves completeness. &
\[
 (Rf_*K)/\varpi^n
   \simeq Rf_*(K_n)
\]
in the bounded constructible range.\\

Exceptional direct image &
Compactify \(f=\bar f j\) and set
\(Rf_!=R\bar f_*j_!\). &
\[
 (Rf_!K)/\varpi^n
   \simeq Rf_!(K_n).
\]\\

Exceptional pullback &
Right adjoint \(Rf^!\) of \(Rf_!\). &
\[
 (Rf^!L)/\varpi^n
   \simeq Rf_n^!(L_n).
\]\\

Tensor &
Derived tensor over \(\widehat{\mathcal O}_X\); on constructible objects it
is already complete. &
\[
 (K\otimes^LL)/\varpi^n
 \simeq K_n\otimes^L_{\Lambda_n}L_n.
\]\\

Internal Hom &
Derived sheaf Hom in the finite-Tor constructible category. &
{\small
\[
 \mathrm{RHom}(K,L)/\varpi^n
 \simeq \mathrm{RHom}(K_n,L_n).
\]
}\\

Duality &
\(\D_XK=R\mathcal Hom(K,\omega_X)\). &
\[
 (\D_XK)/\varpi^n
   \simeq\D_{X,n}(K_n).
\]\\
\bottomrule
\end{longtable}

These formulas are canonical isomorphisms of functors.  They are not
definitional equalities between selected resolutions.  Their coherence under
composition is part of the comparison theorem.

\begin{warning}[The exceptional direct image]
\label{warn:dictionary-fshriek}
Even on the pro-etale site, the constructible \(Rf_!\) is not generally the
right or left derived functor of the degree-zero operation on all abelian
sheaves.  Compactification is part of its definition.  Any formalization
should therefore construct it at the derived level and prove independence of
compactification before exposing the notation \(Rf_!\).
\end{warning}

\section{Perverse structures under the comparison}

For field coefficients \(E\), the equivalence
\eqref{eq:dictionary-rational-equivalence} preserves geometric stalks,
costalks, dimensions of supports, and Verdier duality.  It is therefore
\(t\)-exact for the middle perverse structures:
\[
 \Phi_X({}^p\D^{\leq0}_{\mathrm{adic}})
   ={}^p\D^{\leq0}_{\proet},
\qquad
 \Phi_X({}^p\D^{\geq0}_{\mathrm{adic}})
   ={}^p\D^{\geq0}_{\proet}.
\tag{B.6}\label{eq:dictionary-perverse-texact}
\]
Consequently it commutes canonically with perverse truncation, perverse
cohomology, \(j_{!*}\), and intersection complexes.

Integral coefficients require a two-column dictionary:

\begin{center}
\begin{tabularx}{0.94\textwidth}{>{\raggedright\arraybackslash}p{0.26\textwidth}XX}
\toprule
Operation & \(p\)-heart & \(p^+\)-heart\\
\midrule
Verdier duality &
\(\D:{}^p\Perv(X,\mathcal O)^{\op}\to{}^{p^+}\Perv(X,\mathcal O)\) &
\(\D:{}^{p^+}\Perv(X,\mathcal O)^{\op}\to{}^p\Perv(X,\mathcal O)\)\\
Intermediate extension &
\({}^pj_{!*}\) &
\({}^{p^+}j_{!*}\)\\
Rationalization &
\multicolumn{2}{c}{Both map to the same
\(j_{!*}\) in \(\Perv(X,E)\).}\\
\bottomrule
\end{tabularx}
\end{center}

The two integral middle extensions may differ by boundary
\(\varpi\)-torsion.  Rationalization kills that difference.  The rational
perverse heart is Verdier self-dual and is the category in which weights,
simple objects, pure semisimplicity, and decomposition are used.

\section{Weil structures and Frobenius}

Let \(X_0/\Fq\), let \(X=X_0\times_{\Fq}\Fbar\), and let
\(F_X:X\to X\) be geometric Frobenius in the convention of these notes.  A
Weil complex is
\[
  (K,\phi),\qquad
  \phi:F_X^*K\xrightarrow{\sim}K.
\tag{B.7}\label{eq:dictionary-weil-object}
\]
Under the integral comparison, this is the same as a compatible system of
finite-level Weil structures on \(K_n\).  The comparison isomorphisms for all
six operations commute with \(\phi\), so traces and local terms agree in the
two models.

The Tate object has the three matching presentations
\[
 \Zl(1)=\varprojlim_n\mu_{\ell^n},
 \qquad
 \widehat{\mathcal O}(1)
   =\Zl(1)\widehat{\otimes}_{\Zl}\mathcal O,
\qquad
 E(1)=\widehat{\mathcal O}(1)[\varpi^{-1}],
\qquad
 F\mid E(1)=q^{-1}.
\tag{B.8}\label{eq:dictionary-tate}
\]
The last equality uses geometric Frobenius.  It becomes \(q\) for arithmetic
Frobenius.

\section{What is and is not gained}

\begin{longtable}{
 >{\raggedright\arraybackslash}p{0.44\textwidth}
 >{\raggedright\arraybackslash}p{0.46\textwidth}}
\toprule
Eliminated from the proof spine & Still mathematically required\\
\midrule
\endhead
AR-null objects in every public definition. &
Derived completeness and conservativity of finite reduction.\\

Repeated choices of normalized representatives. &
Canonical comparison isomorphisms and their coherence.\\

Pretending that \(\Ql\) is a sheaf on the small etale site. &
Finite-level etale finiteness and base-change theorems.\\

Defining \(R\Gamma(X,\Ql)\) as an external inverse limit. &
The theorem identifying pro-etale cohomology with continuous etale
cohomology.\\

Coefficient towers in the statements of perverse and Fourier functors. &
Integral \(p/p^+\) bookkeeping and finite coefficient-field descent.\\
\bottomrule
\end{longtable}

Thus the pro-etale model shortens the foundational interface without
weakening it.  The classical model remains available as a comparison target
and as the source of finite-level geometric theorems.
